\documentclass{article}
\usepackage[margin=1in]{geometry}
\usepackage{longtable}
\begin{document}

\section*{Phase 8 Task 4 - Final Metrics, Limitations, and Stop-Condition Review}

\textbf{Source of truth:} \texttt{Documentation/SRS.tex}, the generated metrics-review manifest under \texttt{reports/phase8/metrics\_review/}, and the previously generated Phase 8 Task 2 and Task 3 artifacts.\\
\textbf{Task goal:} Assemble the final metrics bundle in the SRS priority order and write the honest limitations and stop-condition review for the current repo snapshot.

\section*{Bottom-Line Assessment}
The metric framework is now frozen and documented, but it is \textbf{not materially populated} in this workspace. The correct closeout statement for the current snapshot is therefore:
\begin{itemize}
\item the project knows \textbf{which} metrics matter and in what priority order,
\item the repo contains the code paths and policy layers needed to compute them,
\item but the current local evidence base is not strong enough to defend any of the four SRS priority metrics numerically.
\end{itemize}

\section*{Metrics Bundle in SRS Priority Order}
\begin{longtable}{p{0.23\linewidth}p{0.20\linewidth}p{0.45\linewidth}}
\textbf{Metric Family} & \textbf{Current Status} & \textbf{Assessment} \\
Ranked alert precision and usefulness & Not materialized in workspace & There are no persisted alerts, delivery attempts, or analyst-feedback rows in the current local SQLite snapshot, so alert precision/usefulness cannot be computed honestly. \\
Calibration & Not materialized in workspace & No \texttt{model\_evaluation\_runs}, \texttt{calibration\_profiles}, or \texttt{shadow\_model\_scores} rows exist locally. The code supports calibration, but no current evidence demonstrates it. \\
Lead time over public corroboration & Not materialized in workspace & Lead-time analysis needs real alerts and evidence states. Those artifacts are absent in the current workspace. \\
Economic edge under conservative execution & Not materialized in workspace & No \texttt{validation\_runs}, \texttt{backtest\_artifacts}, or \texttt{reports/phase5/} outputs are present locally, so no paper-trade edge claim is defendable. \\
\end{longtable}

\section*{Reference Targets From the SRS}
\begin{itemize}
\item Precision@10 \textgreater{} 0.60
\item Brier score \textless{} 0.20 for calibrated operational probabilities
\item median lead time \textgreater{} 30 minutes on the corroborated subset
\item positive paper-trade edge after fees and slippage
\end{itemize}

These remain \textbf{aspirational targets}, not satisfied outputs in the current workspace snapshot.

\section*{Honest Limitations Review}
\textbf{Weak regimes}
\begin{itemize}
\item The later-phase runtime tables exist structurally, but their local row counts are zero.
\item The current workspace lacks real alert traffic, analyst feedback, model evaluations, and backtest artifacts.
\item The committed Phase 6 trainer is still a linear starter ranker, not the final boosted-tree model described by the full Phase 6 scope.
\end{itemize}

\textbf{Observability caveats}
\begin{itemize}
\item The repo contains a Phase 7 observability / Goodhart study implementation, but there is no local alert/evidence history to populate those metrics here.
\item The observability code itself warns that visible operator precision is not the same thing as preserved lead time or preserved hidden candidate edge.
\end{itemize}

\textbf{Provider issues}
\begin{itemize}
\item The default evidence providers are placeholder \texttt{noop} adapters, so richer real-provider behavior is not the default runtime state.
\item Telegram and Discord integrations exist, but they are environment-gated and may be disabled.
\item The Phase 4 signoff memo explicitly says final validation should use real live candidate-to-alert traffic rather than only local seeded evidence.
\end{itemize}

\textbf{Failure modes}
\begin{itemize}
\item Any current claim of strong alert precision would be overstated because no persisted alert-review evidence exists locally.
\item Any claim of calibrated ML quality would be overstated because thresholds are still framed as shadow-only recommendations.
\item Any claim of full replay-to-alert reproducibility would be overstated because Task 2 marked the reference path as missing runtime outputs.
\end{itemize}

\section*{Stop-Condition Review}
\begin{longtable}{p{0.36\linewidth}p{0.22\linewidth}p{0.28\linewidth}}
\textbf{SRS Stop Condition} & \textbf{Status} & \textbf{Assessment} \\
raw archive gaps exceed 1 hour & Unresolved in current workspace & No local raw archive partitions or manifest rows are present, so archive-gap risk cannot be measured from this snapshot. \\
duplicate trade inflation cannot be controlled & Unresolved in current workspace & The validation framework exists, but no populated local trade evidence proves duplicate inflation is currently bounded. \\
alert false-positive rate remains operationally unusable after suppression tuning & Unresolved in current workspace & No local alert, delivery, or analyst-feedback rows exist, so operational false-positive behavior cannot be assessed. \\
evidence-provider spend exceeds budget without measurable lift & Not triggered in current workspace, but not validated for real providers & Default providers are \texttt{noop} adapters and no live spend evidence exists locally. \\
replay cannot reproduce a historical alert end to end & Active concern & Task 2 froze the reference path but recorded missing runtime outputs across replay, candidate, alert, validation, and ML/research stages. \\
\end{longtable}

\section*{What This Means for Phase 8}
Task 4 should be read as a \textbf{truthfulness checkpoint}. The repo now has a formal metric order, a formal operating mode, and a frozen reference path. What it still does \textbf{not} have in this workspace is the materialized runtime evidence required to defend those metrics numerically. Therefore:
\begin{itemize}
\item the metric framework is ready,
\item the evidence base is not,
\item and the most serious remaining concern is still the lack of a full replay-to-alert end-to-end reproduction packet.
\end{itemize}

\section*{Generated Artifact}
Task 4 also writes:
\begin{itemize}
\item \texttt{reports/phase8/metrics\_review/phase8\_metrics\_review\_manifest.json}
\item \texttt{reports/phase8/metrics\_review/phase8\_metrics\_review\_summary.md}
\end{itemize}

These files record the metric bundle, evidence status, limitations review, and stop-condition ledger in machine-readable form.

\end{document}
